\documentclass[runningheads]{llncs}
\usepackage{cite}
\usepackage{amsmath,amssymb,amsfonts}
\usepackage{algorithmic}
\usepackage{graphicx}
\usepackage{textcomp}
\usepackage{xcolor}
\usepackage{booktabs}
\usepackage{multirow}
\usepackage[hidelinks]{hyperref}
\usepackage{placeins}

\begin{document}

    \title{Towards a Robust and Explainable Pipeline for Diabetic Retinopathy Classification through Quality-Aware GenAI Image Restoration}
\titlerunning{Robust and Explainable DR Pipeline via Quality-Aware GenAI Restoration}

% ANONYMIZED AUTHOR BLOCK FOR DOUBLE-BLIND REVIEW
\author{Anonymous Author(s)}
\authorrunning{Anonymous Author(s)}
\institute{Anonymous Institution(s)}

\maketitle

\begin{abstract}
Diabetic retinopathy (DR) screening models achieve high accuracy on curated datasets, yet $15$--$30\%$ of real-world fundus images suffer from degradations $\mathcal{D}(x, t)$ parameterised by type $t \in \{\text{blur},\, \text{exposure},\, \text{noise}\}$ at three severity levels. A common hypothesis is that generative AI (GenAI) image restoration $\mathcal{R}(\mathcal{D}(x,t),t) \approx x$ can recover diagnostic accuracy on these degraded images. In this paper, we present a controlled robustness benchmark across nine synthetic degradations to evaluate modern Convolutional Neural Networks (CNNs) and Vision Transformers (ViTs). We investigate the effect of restoration on diagnostic performance using six methods, including Generative Adversarial Networks (GANs) and Denoising Diffusion Probabilistic Models (DDPMs). Our findings reveal a counter-intuitive negative result: while GenAI restoration improves pixel fidelity (measured by $\ell_1$ reconstruction error), it does not reliably recover diagnostic accuracy, and diffusion-based restorers actively harm it due to distribution shift $\|p_{\text{restored}} - p_{\text{train}}\|$. Furthermore, we demonstrate that explanation heatmaps $\phi(x)$ degrade faster than accuracy, with Integrated Gradients stability dropping to $\rho < 0.3$ at high degradation. To address these challenges, we propose a quality-aware triage pipeline using selective prediction with a calibrated confidence score $\hat{c} = \max_k\, \sigma(f(x))_k$. At $80\%$ coverage, our triage approach achieves a Quadratic Weighted Kappa ($\kappa_w = 0.665$) on degraded images for the CNN, outperforming both a ``do-nothing'' baseline ($\kappa_w = 0.525$) and a ``restore-all'' strategy ($\kappa_w = 0.464$). Restoration is shown to be beneficial only under severe noise, where $\kappa_w$ rises from $-0.007$ (raw) to $0.433$ with pathology-preserving diffusion.
\keywords{Diabetic Retinopathy \and Image Restoration \and Generative AI \and Explainable AI \and Robustness \and Selective Prediction}
\end{abstract}

%-----------------------------------------------------------------------
\section{Introduction}
%-----------------------------------------------------------------------
Clinical screening for diabetic retinopathy (DR) produces a high volume of fundus photographs. While deep-learning models can grade DR with approximately 85\% accuracy on clean, curated images, real-world screening often yields images of uneven quality. It is estimated that 15--30\% of screening images exhibit motion blur, mis-exposure, or sensor noise. The vulnerability of diagnostic models to such quality degradation is a critical concern, particularly regarding the stability of the predictions and the faithfulness of Explainable AI (XAI) heatmaps.

Recently, generative AI (GenAI) techniques, including Generative Adversarial Networks (GANs) and Denoising Diffusion Probabilistic Models (DDPMs), have become popular for image restoration. A natural hypothesis is that applying these restoration tools to degraded fundus images will recover the lost diagnostic signal. However, it remains unclear whether these models recover the underlying pathology or merely hallucinate plausible but diagnostically irrelevant textures, thereby introducing distribution shifts that harm classification.

In this work, we investigate the following research questions:
\begin{itemize}
    \item \textbf{RQ1:} How does a Vision Transformer (ViT) compare to a modern CNN as image quality degrades?
    \item \textbf{RQ2:} Do explanation methods (e.g., Grad-CAM, Integrated Gradients) remain faithful and stable as quality drops?
    \item \textbf{RQ3:} Can GenAI restoration recover diagnostic accuracy, as opposed to merely improving pixel fidelity?
    \item \textbf{RQ4:} Can a quality-aware system route each image to the optimal pipeline and output a reliable trust score?
\end{itemize}

Our core thesis is that generative restoration improves pixel fidelity but does not reliably recover diagnostic accuracy. In fact, diffusion-based restorers actively harm classification because their output is out-of-distribution (OOD) for the classifier. Instead, a quality-aware triage system that refrain on low-confidence images and applies restoration only where it is measurably beneficial is a safer, more deployable solution.

The contributions of this paper are as follows:
1) A controlled robustness benchmark for DR grading across 9 graded degradations comparing ConvNeXt-Base and CLIP-ViT-B/16.
2) Quantitative evidence that explanations degrade faster than accuracy.
3) A causal isolation of why restoration fails via a clean-image control that measures each restorer's inherent distribution shift.
4) A quantified demonstration that GANs are relatively safe while diffusion models are harmful for downstream diagnosis, highlighting that raw accuracy can hide harm that Quadratic Weighted Kappa (QWK) reveals.
5) A quality-aware triage and selective-prediction system with a calibrated confidence score that outperforms both ``do-nothing'' and ``restore-everything'' approaches for CNNs.

%-----------------------------------------------------------------------
\section{Related Work}
%-----------------------------------------------------------------------
\textbf{Diabetic Retinopathy and Robustness:} DR grading is a well-studied task, often formulated as an ordinal classification problem on datasets like APTOS 2019 \cite{aptos2019}. While models like ResNet \cite{he2016resnet} and ConvNeXt \cite{liu2022convnext} perform well on clean data, their robustness to corruptions varies. Vision Transformers (ViTs) \cite{dosovitskiy2021vit, radford2021clip} have shown distinct robustness profiles compared to CNNs.

\textbf{Image Restoration:} Traditional methods like CLAHE \cite{pizer1987clahe} improve contrast but do not hallucinate details. Modern GenAI approaches include GAN-based super-resolution (e.g., A-ESRGAN \cite{wei2021aesrgan}, SwinIR \cite{liang2021swinir}) and diffusion models (e.g., DDPM \cite{ho2020ddpm}, Cold Diffusion \cite{bansal2022colddiffusion}). Recent evidence suggests diffusion models may underperform at general image restoration in medical contexts \cite{arxiv241209324}.

\textbf{Explainability (XAI):} Methods like Grad-CAM \cite{selvaraju2017gradcam}, SHAP \cite{lundberg2017shap}, Integrated Gradients \cite{sundararajan2017ig}, and Attention Rollout \cite{abnar2020rollout} are used to interpret model decisions. However, the stability of these explanations under distribution shifts remains a challenge.

\textbf{Selective Prediction and Calibration:} Selective classification \cite{geifman2017selective} allows models to avoid from predicting when confidence is low. Proper calibration, such as temperature scaling \cite{guo2017calibration}, is essential for generating reliable confidence scores. Image quality assessment networks like EyeQ \cite{fu2019eyeq} provide prior signals for triage.

%-----------------------------------------------------------------------
\section{Methods}
%-----------------------------------------------------------------------

\subsection{Datasets and Degradation Model}
We utilize the APTOS 2019 Blindness Detection dataset \cite{aptos2019}, which contains fundus images labeled with a 5-grade DR severity scale (0: No DR, 1: Mild, 2: Moderate, 3: Severe, 4: Proliferative). The dataset is heavily class-imbalanced. We also utilize the EyeQ \cite{fu2019eyeq} criteria for per-image quality labels (good, usable, reject).

To create a controlled testbed, we apply synthetic, graded degradations to high-quality images. We implement three degradation primitives: Gaussian blur, exposure (gain) shift, and additive Gaussian noise. Each is applied at three severity levels (low, mid, high), resulting in 9 degradation conditions.

\subsection{Classifiers}
We evaluate two primary base architectures: ConvNeXt-Base (384px) \cite{liu2022convnext} and CLIP-ViT-B/16 (384px) \cite{radford2021clip}. Images undergo Ben Graham preprocessing \cite{graham2015kaggle}. The models employ a multi-scale feature fusion head and are trained using an ordinal focal loss \cite{lin2017focal} to account for the ordinal nature of DR grades. Training incorporates MixUp \cite{zhang2018mixup}, CutMix \cite{yun2019cutmix}, RandAugment \cite{cubuk2020randaugment}, and Exponential Moving Average (EMA). At inference, we apply 4-view Test-Time Augmentation (TTA, averaging softmax over horizontal/vertical flips) uniformly to all models as a standard, cost-free boost; we note that it disproportionately helps under severe noise, where single-view predictions are high-variance, and we therefore report the restoration comparisons (Table~\ref{tab:qwk_recovery}) single-view so that raw and restored images are evaluated on equal footing. Note that an EfficientNetV2-S model collapsed during training and is excluded from the primary robustness claims.

\subsection{Restoration Models}
We evaluate six restoration strategies:
1) \textbf{CLAHE} \cite{pizer1987clahe} as a non-generative baseline.
2) \textbf{A-ESRGAN} \cite{wei2021aesrgan}, a GAN-based super-resolution model.
3) \textbf{SwinIR+GAN} \cite{liang2021swinir}.
4) \textbf{Cold Diffusion} \cite{bansal2022colddiffusion}, which uses deterministic degradation as the forward process.
5) \textbf{Vanilla DDPM} \cite{ho2020ddpm} (conditional).
6) \textbf{Pathology-preserving DDPM} (a custom variant that heavily guides the reverse process using the structural edges of the input image to prevent lesion hallucination).

\subsection{Explainability (XAI) Methods}
We apply Grad-CAM \cite{selvaraju2017gradcam} for CNNs, Attention Rollout \cite{abnar2020rollout} for ViTs, and Integrated Gradients \cite{sundararajan2017ig} across architectures. We omit SHAP as its perturbation-based sampling is computationally prohibitive across our large degradation grid and introduces stochastic noise that would confound the stability measurement; IG provides a deterministic, architecture-agnostic attribution suitable for cross-model comparison.XAI stability is measured via the Spearman rank correlation of heatmaps under degradation compared to the clean baseline.

\subsection{Metrics and Calibration}
Due to class imbalance, the headline metric is Quadratic Weighted Kappa (QWK) \cite{cohen1968kappa}, supplemented by accuracy. We report bootstrap 95\% confidence intervals (CIs) to ensure rigor given the test set size. Model calibration is performed using temperature scaling, measured by Expected Calibration Error (ECE) \cite{guo2017calibration}.

\subsection{Quality-Aware Triage and Selective Prediction}
We implement a triage system using a quality classifier to assign confidence scores based on the calibrated softmax maximum. The protocol utilizes selective prediction: images below a confidence threshold are rejected (flagged for re-acquisition). For the remaining images, a fixed prior routing rule is applied: images with severe noise are restored using A-ESRGAN, while all others are passed raw to the classifier. We specifically select A-ESRGAN over DDPM-pathology for noise restoration, despite the latter's slightly higher peak QWK, because our clean-image control demonstrates that A-ESRGAN is distribution-safe (+0.025 accuracy on clean images) whereas the DDPM introduces harmful shift ($-0.125$ accuracy).

%-----------------------------------------------------------------------
\section{Experiments and Results}
%-----------------------------------------------------------------------

\subsection{RQ1: Model Robustness Benchmark}
Table~\ref{tab:rq1_robustness} presents our controlled robustness benchmark across the 9 degradation conditions. The results indicate that the CLIP-ViT-B/16 degrades more gracefully than the ConvNeXt-Base under severe blur and noise---a gap that is invisible on clean benchmarks. For instance, under severe noise, the ViT retains a QWK of 0.522 compared to the ConvNeXt-Base's 0.415. It is important to note that the ViT's robustness may stem from its large-scale contrastive pretraining (CLIP) rather than solely its architecture. Interestingly, low-severity exposure and noise perturbations slightly outperform the clean baseline for both models, likely acting as a regularizing test-time augmentation.

\begin{table}[!t]
\caption{RQ1: Robustness Benchmark (Canonical QWK under Degradation, 4-view TTA). These canonical numbers use test-time augmentation; the single-view Raw values in Table~\ref{tab:qwk_recovery} are therefore lower.}
\label{tab:rq1_robustness}
\begin{center}
\begin{tabular}{lcc}
\toprule
\textbf{Condition} & \textbf{ConvNeXt-Base} & \textbf{CLIP-ViT-B/16} \\
\midrule
Clean (Baseline) & 0.625 & \textbf{0.707} \\
\midrule
Blur Low & \textbf{0.868} & 0.807 \\
Blur Mid & 0.515 & \textbf{0.648} \\
Blur High & 0.429 & \textbf{0.550} \\
\midrule
Exposure Low & \textbf{0.929} & 0.899 \\
Exposure Mid & \textbf{0.911} & 0.879 \\
Exposure High & \textbf{0.838} & 0.823 \\
\midrule
Noise Low & \textbf{0.889} & 0.876 \\
Noise Mid & 0.651 & \textbf{0.684} \\
Noise High & 0.415 & \textbf{0.522} \\
\bottomrule
\end{tabular}
\end{center}
\end{table}

\subsection{RQ2: XAI Stability}
Explanation stability falls rapidly under degradation, even when the model's accuracy remains reasonable. Table~\ref{tab:rq2_xai} shows the stability (Spearman correlation) of Integrated Gradients (IG) heatmaps compared to the clean baseline. The scores drop significantly under high degradation. This demonstrates that explanations drift faster than accuracy, indicating that models may be ``right for the wrong reasons'' under degradation.

\begin{table}[!t]
\caption{RQ2: XAI Stability (Integrated Gradients Spearman Correlation)}
\label{tab:rq2_xai}
\begin{center}
\begin{tabular}{lccc}
\toprule
\textbf{Model} & \textbf{Blur High} & \textbf{Exposure High} & \textbf{Noise High} \\
\midrule
ConvNeXt-Base & 0.117 & 0.210 & 0.128 \\
CLIP-ViT-B/16 & \textbf{0.184} & \textbf{0.302} & \textbf{0.224} \\
\bottomrule
\end{tabular}
\end{center}
\end{table}

\subsection{RQ3: Generative Restoration and Distribution Shift}
To isolate the cause of restoration failure, we performed a clean-image control experiment (Table~\ref{tab:dist_shift}). Clean, un-degraded images were passed through each restorer and re-classified using the ConvNeXt-Base model.

\begin{table}[!t]
\caption{Distribution-Shift Control (Clean $\rightarrow$ Restore $\rightarrow$ Re-classify, ConvNeXt-Base; $n=40$). Baseline clean accuracy = 0.85.}
\label{tab:dist_shift}
\begin{center}
\begin{tabular}{lcc}
\toprule
\textbf{Restorer} & \textbf{Accuracy} & \textbf{$\Delta$ vs clean} \\
\midrule
\textbf{A-ESRGAN (GAN)} & \textbf{0.875} & \textbf{$+0.025$} \\
CLAHE & 0.775 & $-0.075$ \\
DDPM-pathology & 0.725 & $-0.125$ \\
SwinIR+GAN & 0.675 & $-0.175$ \\
Cold Diffusion & 0.300 & $-0.550$ \\
DDPM (vanilla) & 0.100 & $-0.750$ \\
\bottomrule
\end{tabular}
\end{center}
\end{table}

The results show that A-ESRGAN is the only distribution-safe restorer: its change relative to baseline ($+0.025$) is within the sampling noise of this $n=40$ control, indicating it preserves rather than improves the classifier's input distribution. Diffusion variants introduce massive distribution shifts that destroy classification performance, even on pristine images. Notably, while the vanilla DDPM collapsed due to under-training, Cold Diffusion was properly trained and still exhibited a massive accuracy drop (0.300), isolating the inherent harm of the diffusion process.

Table~\ref{tab:qwk_recovery} presents the QWK recovery on the degraded test set for the ConvNeXt-Base model. The raw (do-nothing) approach wins on blur and exposure. Restoration substantially helps \textit{only} at severe noise, where A-ESRGAN (0.336) and DDPM-pathology (0.433) recover performance from near zero.

\begin{table}[!t]
\caption{QWK by Restorer vs Raw (Degraded Test Set, ConvNeXt-Base). Bold = best per row. All pipelines are evaluated single-view (no TTA) so that Raw and restored images are compared on equal footing; hence Raw QWK here is lower than the TTA values in Table~\ref{tab:rq1_robustness}.}
\label{tab:qwk_recovery}
\begin{center}
\setlength{\tabcolsep}{3.5pt}
\small
\begin{tabular}{lrrrrrrr}
\toprule
\textbf{Deg./Level} & \textbf{Raw} & \textbf{CLAHE} & \textbf{A-ESR} & \textbf{SwinIR} & \textbf{C-Diff} & \textbf{DDPM} & \textbf{D-path} \\
\midrule
Blur Low      & \textbf{.650} & .591 & .602 & .539 & .581 & .031 & .313 \\
Blur Mid      & \textbf{.449} & .445 & .238 & .408 & .301 & $-.004$ & .206 \\
Blur High     & .239 & \textbf{.251} & .172 & .200 & .091 & .052 & .061 \\
Exp. Low      & \textbf{.807} & .700 & .633 & .482 & .614 & $-.025$ & .472 \\
Exp. Mid      & \textbf{.744} & .629 & .496 & .651 & .620 & .000 & .455 \\
Exp. High     & .509 & \textbf{.590} & .359 & .466 & .503 & .011 & .367 \\
Noise Low     & \textbf{.813} & .746 & .625 & .560 & .536 & .021 & .408 \\
Noise Mid     & \textbf{.674} & .483 & .519 & .564 & $-.005$ & $-.030$ & .295 \\
\textbf{Noise High} & $-.007$ & $-.035$ & .336 & .065 & .001 & .023 & \textbf{.433} \\
\bottomrule
\end{tabular}
\end{center}
\end{table}

\subsection{RQ4: Quality-Aware Triage and Selective Prediction}
We evaluate our triage pipeline against ``do-nothing'' and ``restore-all'' baselines (Table~\ref{tab:triage}).

\begin{table}[!t]
\caption{Selective Prediction / Triage at Full and 80\% Coverage (Degraded Stream)}
\label{tab:triage}
\begin{center}
\begin{tabular}{llcc|cc}
\toprule
 & & \multicolumn{2}{c|}{\textbf{Full Coverage (100\%)}} & \multicolumn{2}{c}{\textbf{Selective (80\%)}} \\
\textbf{Model} & \textbf{Pipeline} & \textbf{Acc.} & \textbf{QWK} & \textbf{Acc.} & \textbf{QWK} \\
\midrule
ConvNeXt-Base & Do-nothing  & 0.619 & 0.465 & 0.691 & 0.525 \\
ConvNeXt-Base & Restore-all & 0.602 & 0.442 & 0.674 & 0.464 \\
ConvNeXt-Base & \textbf{Triage}    & \textbf{0.656} & \textbf{0.582} & \textbf{0.739} & \textbf{0.665} \\
\midrule
ViT-B/16 & Do-nothing  & \textbf{0.638} & \textbf{0.621} & \textbf{0.706} & \textbf{0.701} \\
ViT-B/16 & Restore-all & 0.561 & 0.388 & 0.621 & 0.441 \\
ViT-B/16 & Triage      & 0.636 & 0.596 & 0.698 & 0.656 \\
\bottomrule
\end{tabular}
\end{center}
\end{table}

For the ConvNeXt-Base model, the triage pipeline achieves 0.739 accuracy and 0.665 QWK at 80\% coverage, significantly outperforming do-nothing (0.525 QWK) and restore-all (0.464 QWK). Notably, the ``restore-all'' strategy maintains acceptable accuracy but causes QWK to crater, as indiscriminate restoration destroys the ordinal structure of minority classes. For the already-robust ViT, abstention is the primary lever, and triage performs similarly to (and slightly worse than) do-nothing, indicating that triage benefits weaker CNN models the most.

\subsection{Calibration}
Proper calibration enables the selective prediction threshold. After temperature scaling, the ViT's ECE improved from 0.071 to 0.051, validating the use of its softmax outputs for confidence scoring. Conversely, the ConvNeXt-Base was already well-calibrated (ECE remained stable at 0.070, Temperature=1.000), meaning its triage gains were driven by its native confidence distributions rather than post-hoc scaling.

\FloatBarrier

%-----------------------------------------------------------------------
\section{Discussion}
%-----------------------------------------------------------------------
Our findings highlight that visual fidelity does not equate to diagnostic accuracy. Generative restoration introduces severe out-of-distribution shifts. GAN-based super-resolution (A-ESRGAN) proved to be the only distribution-safe restorer in our tests. The discrepancy between accuracy and QWK in the ``restore-all'' setting underscores the danger of using accuracy alone for imbalanced medical datasets. Clinically, our results suggest that ungradable images should be flagged for re-acquisition rather than hallucinated via restoration, with restoration reserved exclusively for specific, severe noise profiles.

%-----------------------------------------------------------------------
\section{Limitations}
%-----------------------------------------------------------------------
This study relies on synthetic degradations applied to a single dataset (APTOS). Real-world degradations are more complex. The vanilla DDPM model collapsed to noise, representing a limitation in our diffusion evaluation, though Cold Diffusion provides a properly trained alternative point of evidence. Finally, the CLIP-ViT's robustness advantage may stem from its large-scale pretraining rather than its architecture alone.

%-----------------------------------------------------------------------
\section{Future Work}
%-----------------------------------------------------------------------
Future research should explore diffusion-based restoration using paired degraded/clean training data \cite{arxiv230809388}. Leveraging retinal foundation models like RETFound \cite{zhou2023retfound} and applying test-time adaptation \cite{wang2021tent} could provide stronger robustness guarantees. Validating the triage pipeline on datasets with natural, non-synthetic degradations is a critical next step.

%-----------------------------------------------------------------------
\section{Conclusion}
%-----------------------------------------------------------------------
We demonstrated that generative image restoration is not a panacea for degraded diabetic retinopathy screening images. Diffusion models actively harm classification due to distribution shifts, and restoring all images destroys ordinal predictive performance. Instead, a quality-aware triage system that abstains on low-confidence images and selectively applies GAN-based restoration only to severely noisy images provides a robust, clinically viable alternative for CNN-based pipelines.

\bibliographystyle{splncs04}
\bibliography{references}

\end{document}
